%=============================================================
% Metodi Numerici per Data Science
% Appunti del Corso
%=============================================================
\documentclass[12pt, a4paper, twoside]{book}

%------------------------------------------------------------
% Encoding e lingua
%------------------------------------------------------------
\usepackage[utf8]{inputenc}
\usepackage[T1]{fontenc}
\usepackage[italian]{babel}

%------------------------------------------------------------
% Matematica
%------------------------------------------------------------
\usepackage{amsmath}
\usepackage{amssymb}
\usepackage{amsthm}
\usepackage{mathtools}

%------------------------------------------------------------
% Layout e tipografia
%------------------------------------------------------------
\usepackage[
  top=2.5cm,
  bottom=2.5cm,
  inner=3cm,
  outer=2cm
]{geometry}
\usepackage{microtype}
\usepackage{setspace}
\onehalfspacing

%------------------------------------------------------------
% Intestazioni e piè di pagina
%------------------------------------------------------------
\usepackage{fancyhdr}
\pagestyle{fancy}
\fancyhf{}
\fancyhead[LE,RO]{\thepage}
\fancyhead[LO]{\nouppercase{\rightmark}}
\fancyhead[RE]{\nouppercase{\leftmark}}
\renewcommand{\headrulewidth}{0.4pt}

%------------------------------------------------------------
% Colori e grafica
%------------------------------------------------------------
\usepackage[dvipsnames]{xcolor}
\usepackage{graphicx}
\graphicspath{{images/}}

%------------------------------------------------------------
% Tabelle e liste
%------------------------------------------------------------
\usepackage{booktabs}
\usepackage{array}
\usepackage{enumitem}

%------------------------------------------------------------
% Algoritmi
%------------------------------------------------------------
\usepackage{algorithm}
\usepackage{algpseudocode}
\algnewcommand\algorithmicforeach{\textbf{for each}}
\algdef{S}[FOR]{ForEach}[1]{\algorithmicforeach\ #1\ \textbf{do}}

%------------------------------------------------------------
% Referenze e link
%------------------------------------------------------------
\usepackage[
  colorlinks=true,
  linkcolor=NavyBlue,
  citecolor=ForestGreen,
  urlcolor=RoyalBlue
]{hyperref}
\usepackage[italian]{cleveref}

%------------------------------------------------------------
% Ambienti teoremi
%------------------------------------------------------------
\theoremstyle{plain}
\newtheorem{theorem}{Teorema}[chapter]
\newtheorem{lemma}[theorem]{Lemma}
\newtheorem{corollary}[theorem]{Corollario}
\newtheorem{proposition}[theorem]{Proposizione}

\theoremstyle{definition}
\newtheorem{definition}[theorem]{Definizione}
\newtheorem{example}[theorem]{Esempio}
\newtheorem{remark}[theorem]{Osservazione}

%------------------------------------------------------------
% Notazioni
%------------------------------------------------------------
\newcommand{\R}{\mathbb{R}}
\newcommand{\N}{\mathbb{N}}
\newcommand{\F}{\mathbb{F}}
\newcommand{\norm}[1]{\left\|#1\right\|}
\newcommand{\abs}[1]{\left|#1\right|}
\newcommand{\fl}{\operatorname{fl}}
\newcommand{\eps}{\varepsilon}

%=============================================================
\begin{document}

%------------------------------------------------------------
% Frontespizio
%------------------------------------------------------------
\begin{titlepage}
  \centering
  \vspace*{2cm}

  {\Large\scshape Università degli Studi\par}
  \vspace{0.5cm}
  {\large Corso di Laurea in Data Science\par}

  \vspace{3cm}

  {\Huge\bfseries Metodi Numerici\par}
  {\Huge\bfseries per Data Science\par}

  \vspace{0.5cm}
  \rule{0.6\textwidth}{0.4pt}

  \vspace{1cm}
  {\large\itshape Appunti del corso\par}

  \vfill

  \begin{minipage}{0.45\textwidth}
    \begin{flushleft}
      \textbf{Studente}\\
      Edoardo Laghi
    \end{flushleft}
  \end{minipage}
  \hfill
  \begin{minipage}{0.45\textwidth}
    \begin{flushright}
      \textbf{Anno Accademico}\\
      2025/2026
    \end{flushright}
  \end{minipage}

  \vspace{2cm}
\end{titlepage}

%------------------------------------------------------------
% Pagina dei diritti / note
%------------------------------------------------------------
\thispagestyle{empty}
\null\vfill
\noindent
\textit{Questi appunti sono stati redatti durante il corso di Metodi Numerici per Data
Science nell'Anno Accademico 2025/2026. Il materiale è basato sulle lezioni tenute a
lezione e rappresenta una rielaborazione personale del programma del corso.}

\vspace{1cm}
\noindent
\textit{Ultima revisione: \today}
\newpage

%------------------------------------------------------------
% Indice
%------------------------------------------------------------
\tableofcontents
\newpage

%------------------------------------------------------------
% Prefazione
%------------------------------------------------------------
\chapter*{Prefazione}
\addcontentsline{toc}{chapter}{Prefazione}

Questi appunti raccolgono i contenuti delle lezioni del corso di \emph{Metodi Numerici
per Data Science}. Il corso si propone di introdurre i principali strumenti dell'analisi
numerica con applicazioni alla scienza dei dati: dall'aritmetica in virgola mobile, alla
risoluzione di sistemi lineari, all'approssimazione di funzioni e alla decomposizione
spettrale di matrici.

\medskip
L'obiettivo è fornire una comprensione teorica dei metodi, insieme a indicazioni
pratiche sulla loro implementazione ed efficienza computazionale.

\medskip
Il materiale è organizzato come segue:
\begin{itemize}
  \item \textbf{\Cref{ch:fp}}: Aritmetica in virgola mobile, standard IEEE 754,
        epsilon di macchina e propagazione degli errori.
  \item \textbf{\Cref{ch:sistemi}}: Sistemi lineari, spazi vettoriali, norme e numero
        di condizionamento.
  \item \textbf{\Cref{ch:diretti}}: Metodi diretti per la risoluzione di sistemi
        lineari (eliminazione di Gauss, fattorizzazione LU).
  \item \textbf{\Cref{ch:iterativi}}: Metodi iterativi stazionari (Jacobi,
        Gauss-Seidel) e criteri di convergenza.
  \item \textbf{\Cref{ch:condizionamento}}: Analisi del condizionamento dei sistemi
        lineari e perturbazioni.
  \item \textbf{\Cref{ch:interpolazione}}: Interpolazione polinomiale, polinomio di
        Lagrange, errore di interpolazione, costante di Lebesgue.
  \item \textbf{\Cref{ch:minimi_quadrati}}: Problema dei minimi quadrati,
        fattorizzazione QR e matrici di Householder.
  \item \textbf{\Cref{ch:autovalori}}: Autovalori, cerchi di Gershgorin,
        decomposizione ai valori singolari (SVD) e algoritmo QR.
\end{itemize}

%------------------------------------------------------------
% Capitoli
%------------------------------------------------------------
%=============================================================
% Capitolo 1 — Aritmetica in Virgola Mobile
%=============================================================
\chapter{Aritmetica in Virgola Mobile}
\label{ch:fp}

\section{Introduzione: il problema della rappresentazione}

Nei calcolatori moderni i numeri reali vengono rappresentati con un numero
\emph{finito} di cifre binarie. Questo comporta due conseguenze fondamentali:
\begin{itemize}
  \item non tutti i numeri reali sono rappresentabili esattamente;
  \item le operazioni aritmetiche introducono errori di arrotondamento.
\end{itemize}

Comprendere questi limiti è essenziale per valutare l'affidabilità dei risultati
numerici.

%------------------------------------------------------------
\section{Sistema posizionale e notazione scientifica}
%------------------------------------------------------------

Un numero reale $x$ viene espresso in \emph{notazione scientifica normalizzata} in
base $b$ come:
\begin{equation}
  x = \pm\, m \cdot b^e,
  \label{eq:scientific}
\end{equation}
dove:
\begin{itemize}
  \item $b \geq 2$ è la \emph{base} del sistema (tipicamente $b = 2$ per i calcolatori);
  \item $m$ è la \emph{mantissa} (o significando), con $1 \leq m < b$;
  \item $e$ è l'\emph{esponente} (intero con segno).
\end{itemize}

\begin{example}[Notazione scientifica in base 10]
Il numero $\pi \approx 3{,}141592\ldots$ si scrive come
\[
  3{,}141592 \times 10^{0}.
\]
Con 5 cifre disponibili per la mantissa, la rappresentazione è
$3{,}14159 \times 10^{0}$ (troncamento) oppure $3{,}14159 \times 10^{0}$ (arrotondamento).
\end{example}

%------------------------------------------------------------
\section{Lo standard IEEE 754}
%------------------------------------------------------------

Lo standard \textbf{IEEE 754} (1985, aggiornato nel 2008 e 2019) definisce il formato
dei numeri in virgola mobile usati in praticamente tutti i processori moderni.

\subsection{Formato della parola binaria}

Un numero floating-point è memorizzato in una parola di $w$ bit suddivisa in tre campi:

\begin{center}
\begin{tabular}{|c|c|c|}
\hline
\textbf{Segno} (1 bit) & \textbf{Esponente} ($e$ bit) & \textbf{Mantissa} ($m$ bit) \\
\hline
\end{tabular}
\end{center}

\medskip
I due formati principali sono riassunti in \cref{tab:ieee754}.

\begin{table}[htbp]
  \centering
  \caption{Formati IEEE 754 a precisione singola e doppia.}
  \label{tab:ieee754}
  \begin{tabular}{lcccc}
    \toprule
    \textbf{Formato} & \textbf{Bit totali} & \textbf{Bit segno} &
    \textbf{Bit esponente} & \textbf{Bit mantissa} \\
    \midrule
    Precisione singola & 32 & 1 & 8  & 23 \\
    Precisione doppia  & 64 & 1 & 11 & 52 \\
    \bottomrule
  \end{tabular}
\end{table}

\subsection{Il bit nascosto}

Nella notazione normalizzata la prima cifra della mantissa è sempre 1 (in base 2),
quindi non viene memorizzata esplicitamente: si guadagna così un bit di precisione
\emph{gratuitamente}. Questo bit è detto \textbf{bit nascosto} (o \emph{implicit
leading bit}).

La mantissa effettiva è quindi $1.f_1f_2\ldots f_m$ dove $f_i \in \{0,1\}$.

\subsection{Configurazioni speciali}

Con $e_\text{bit}$ bit per l'esponente si hanno $2^{e_\text{bit}}$ configurazioni
possibili. Due di queste sono riservate a valori speciali:
\begin{itemize}
  \item \texttt{+Inf}, \texttt{-Inf} — overflow;
  \item \texttt{NaN} (Not a Number) — operazione non definita (es.\ $0/0$).
\end{itemize}
Le rimanenti $2^{e_\text{bit}} - 2$ codificano gli esponenti normali. Ad esempio in
precisione singola si hanno $256 - 2 = 254$ esponenti utilizzabili.

\subsection{Numeri denormalizzati}

I \textbf{numeri denormalizzati} (o subnormali) permettono di rappresentare valori
molto piccoli, al di sotto del minimo normalizzato, a prezzo di una riduzione della
precisione. In questo caso il bit nascosto viene posto a 0:
\[
  x = \pm\, 0.f_1f_2\ldots f_m \times b^{e_\min}.
\]

%------------------------------------------------------------
\section{Insieme dei numeri di macchina}
%------------------------------------------------------------

\begin{definition}[Insieme dei numeri di macchina]
L'insieme $\F \subset \R$ dei \emph{numeri di macchina} in base $b$, con $t$ cifre di
mantissa ed esponente nell'intervallo $[e_\min, e_\max]$, è:
\begin{equation}
  \F(b, t, e_\min, e_\max) = \bigl\{ x \in \R \;:\; x = \pm\, b^e \sum_{i=1}^{t} d_i\, b^{-i},\;
  e \in [e_\min, e_\max],\; d_i \in \{0,\ldots,b{-}1\},\; d_1 \neq 0 \bigr\} \cup \{0\}.
  \label{eq:macchina}
\end{equation}
\end{definition}

\begin{remark}
$\F$ è un insieme \emph{finito} e \emph{discreto}: i numeri di macchina non sono
distribuiti uniformemente sulla retta reale, ma più densamente vicino allo zero e più
rade man mano che ci si allontana.
\end{remark}

%------------------------------------------------------------
\section{Troncamento e arrotondamento}
%------------------------------------------------------------

Sia $x \in \R$ un numero non rappresentabile esattamente in $\F$. L'operazione
$\fl : \R \to \F$ che associa a $x$ il numero di macchina più vicino si chiama
\textbf{arrotondamento}.

\subsection{Troncamento}

Il \emph{troncamento} della mantissa a $t$ cifre è:
\begin{equation}
  \fl_\text{tronc}(x) = \text{sign}(x) \cdot b^e \cdot \sum_{i=1}^{t} d_i\, b^{-i}.
\end{equation}

\subsection{Arrotondamento al più vicino}

L'\emph{arrotondamento al più vicino} (default IEEE) sceglie il numero di macchina più
prossimo a $x$:
\begin{equation}
  \fl_\text{arr}(x) = \arg\min_{y \in \F} |x - y|.
\end{equation}

In caso di equidistanza si applica la regola \emph{round half to even} (arrotondamento
bancario).

%------------------------------------------------------------
\section{Errore di arrotondamento ed epsilon di macchina}
%------------------------------------------------------------

\begin{definition}[Errore relativo di rappresentazione]
L'errore relativo commesso nell'approssimare $x \neq 0$ con $\fl(x)$ è:
\begin{equation}
  \delta = \frac{\fl(x) - x}{x}.
  \label{eq:errore_rel}
\end{equation}
\end{definition}

\begin{definition}[Epsilon di macchina]
L'\emph{epsilon di macchina} $\eps_M$ è il più piccolo numero positivo tale che
\begin{equation}
  \fl(1 + \eps_M) > 1.
  \label{eq:eps}
\end{equation}
Equivalentemente, $\eps_M$ è la distanza relativa tra 1 e il successivo numero di
macchina.
\end{definition}

Per il formato IEEE 754 si ha:
\begin{align}
  \eps_M^{(32)} &= 2^{-23} \approx 1.19 \times 10^{-7} \quad \text{(singola precisione)}, \\
  \eps_M^{(64)} &= 2^{-52} \approx 2.22 \times 10^{-16} \quad \text{(doppia precisione)}.
\end{align}

\begin{theorem}[Limitazione dell'errore relativo]
\label{thm:errore_rel}
Per ogni $x \in \R$ con $|x|$ nel range dei numeri normalizzati vale:
\begin{equation}
  |\delta| = \left|\frac{\fl(x) - x}{x}\right| \leq \frac{1}{2}\, b^{1-t} = \frac{\eps_M}{2}.
  \label{eq:bound_errore}
\end{equation}
Equivalentemente, $\fl(x) = x(1 + \delta)$ con $|\delta| \leq \eps_M/2$.
\end{theorem}

\begin{remark}[Underflow]
Se un numero $x \neq 0$ è troppo piccolo in valore assoluto per essere rappresentato
come normalizzato (e anche come denormalizzato), si ha \emph{underflow}: il numero
viene arrotondato a zero. Questo non è la stessa cosa di $x = 0$.
\end{remark}

%------------------------------------------------------------
\section{Algoritmo di Horner}
\label{sec:horner}
%------------------------------------------------------------

La valutazione di un polinomio
\begin{equation}
  p(x) = a_n x^n + a_{n-1} x^{n-1} + \cdots + a_1 x + a_0
  \label{eq:polinomio}
\end{equation}
richiede, con la formula diretta, $n$ addizioni e $\frac{n(n+1)}{2}$ moltiplicazioni.

L'\textbf{algoritmo di Horner} (o schema di Ruffini–Horner) riduce il costo a $n$
addizioni e $n$ moltiplicazioni, riscrivendo il polinomio in forma annidiata:
\begin{equation}
  p(x) = a_0 + x\bigl(a_1 + x\bigl(a_2 + \cdots + x(a_{n-1} + x\, a_n)\cdots\bigr)\bigr).
  \label{eq:horner}
\end{equation}

\begin{algorithm}[htbp]
  \caption{Algoritmo di Horner per la valutazione di polinomi}
  \label{alg:horner}
  \begin{algorithmic}[1]
    \Require Coefficienti $a_0, a_1, \ldots, a_n$; punto di valutazione $x$
    \Ensure $p \approx p(x)$
    \State $p \gets a_n$
    \For{$k = n-1$ \textbf{downto} $0$}
      \State $p \gets p \cdot x + a_k$
    \EndFor
    \State \Return $p$
  \end{algorithmic}
\end{algorithm}

Il costo computazionale dell'Algoritmo~\ref{alg:horner} è $\mathcal{O}(n)$.

%------------------------------------------------------------
\section{Propagazione degli errori}
\label{sec:propagazione}
%------------------------------------------------------------

Sia $f : \R^n \to \R$ una funzione differenziabile. Se i dati $x_1,\ldots,x_n$ sono
affetti da errori assoluti $\Delta x_1,\ldots,\Delta x_n$, l'errore assoluto sul
risultato è approssimato dalla formula di propagazione:
\begin{equation}
  \Delta f \approx \sum_{i=1}^{n} \left|\frac{\partial f}{\partial x_i}\right| \Delta x_i.
  \label{eq:propagazione}
\end{equation}

In termini di errori relativi (con $\Delta x_i = |\eps_i| |x_i|$):
\begin{equation}
  \frac{\Delta f}{|f|} \approx \sum_{i=1}^{n} \left|\frac{\partial f}{\partial x_i}\right|
  \frac{|x_i|}{|f|} |\eps_i|.
  \label{eq:propagazione_rel}
\end{equation}

\begin{example}[Sottrazione catastrofica]
La sottrazione di due numeri quasi uguali $a \approx b$ amplifica enormemente l'errore
relativo. Se $a = 1{,}0001$ e $b = 1{,}0000$, entrambi con errore relativo $10^{-4}$,
l'errore relativo sulla differenza $a - b = 0{,}0001$ può raggiungere il 100\%.
Questo fenomeno è noto come \emph{cancellazione numerica}.
\end{example}

%------------------------------------------------------------
\section{Aritmetica in virgola mobile IEEE}
\label{sec:aritmetica_ieee}
%------------------------------------------------------------

Lo standard IEEE 754 garantisce che le operazioni elementari $+,-,\times,\div$ (e la
radice quadrata) siano eseguite come se il calcolo fosse fatto con precisione infinita
e poi arrotondato al più vicino numero di macchina:
\begin{equation}
  \fl(x \circ y) = \fl(x \circ_\text{esatta} y) = (x \circ y)(1 + \delta),
  \quad |\delta| \leq \frac{\eps_M}{2},
  \label{eq:ieee_op}
\end{equation}
dove $\circ \in \{+,-,\times,\div\}$.

\begin{remark}[Registri estesi]
Alcuni processori usano registri interni a 80 bit per i calcoli intermedi, riducendo
l'accumulo degli errori di arrotondamento.
\end{remark}

%=============================================================
% Capitolo 2 — Sistemi Lineari
%=============================================================
\chapter{Sistemi Lineari}
\label{ch:sistemi}

\section{Introduzione}

Molti problemi applicativi in data science, ottimizzazione e simulazione si riducono
alla risoluzione di un \textbf{sistema lineare}:
\begin{equation}
  Ax = b,
  \label{eq:sistema}
\end{equation}
dove $A \in \R^{n \times n}$ è la matrice dei coefficienti, $b \in \R^n$ è il termine
noto e $x \in \R^n$ è il vettore incognito da determinare.

%------------------------------------------------------------
\section{Richiami di algebra lineare}
%------------------------------------------------------------

\subsection{Matrici e operazioni fondamentali}

Una matrice $A \in \R^{m \times n}$ ha $m$ righe e $n$ colonne. L'elemento in
posizione $(i,j)$ si indica con $a_{ij}$.

Forma esplicita:
\begin{equation}
  A =
  \begin{pmatrix}
    a_{11} & a_{12} & \cdots & a_{1n} \\
    a_{21} & a_{22} & \cdots & a_{2n} \\
    \vdots & \vdots & \ddots & \vdots \\
    a_{m1} & a_{m2} & \cdots & a_{mn}
  \end{pmatrix}.
\end{equation}

La \textbf{matrice trasposta} $A^\top \in \R^{n \times m}$ ha elementi $(A^\top)_{ij} = a_{ji}$.
La \textbf{matrice coniugata trasposta} (o aggiunta hermitiana) è $A^* = \bar{A}^\top$.

\subsection{Range e rango}

\begin{definition}[Range e nucleo]
Sia $A \in \R^{m \times n}$.
\begin{itemize}
  \item Il \emph{range} (o immagine) di $A$ è $\mathcal{R}(A) = \{Ax : x \in \R^n\} \subseteq \R^m$.
  \item Il \emph{nucleo} di $A$ è $\mathcal{N}(A) = \{x \in \R^n : Ax = 0\} \subseteq \R^n$.
\end{itemize}
\end{definition}

\begin{definition}[Rango]
Il \emph{rango} di $A$, indicato $\operatorname{rank}(A)$, è la dimensione di
$\mathcal{R}(A)$.
\end{definition}

\subsection{Determinante}

Per una matrice quadrata $A \in \R^{n \times n}$:
\begin{itemize}
  \item $\det(A) \neq 0 \iff A$ è non singolare (invertibile);
  \item Se $A$ è triangolare, $\det(A) = \prod_{i=1}^n a_{ii}$;
  \item $\det(AB) = \det(A)\det(B)$.
\end{itemize}

\subsection{Matrice inversa}

Se $A$ è non singolare, esiste unica $A^{-1}$ tale che $AA^{-1} = A^{-1}A = I$.

\begin{remark}
Calcolare $A^{-1}$ esplicitamente è \emph{costoso} ($\mathcal{O}(n^3)$) e spesso
\emph{mal condizionato}: in pratica si risolve $Ax = b$ senza mai calcolare $A^{-1}$.
\end{remark}

\subsection{Autovalori e autovettori}

\begin{definition}
$\lambda \in \mathbb{C}$ è un \emph{autovalore} di $A$ e $v \neq 0$ il corrispondente
\emph{autovettore} se
\begin{equation}
  Av = \lambda v.
  \label{eq:autovalore}
\end{equation}
\end{definition}

Gli autovalori sono le radici del \emph{polinomio caratteristico} $\det(A - \lambda I) = 0$.

\begin{definition}[Raggio spettrale]
Il \emph{raggio spettrale} di $A$ è:
\begin{equation}
  \rho(A) = \max_i |\lambda_i|,
  \label{eq:raggio_spettrale}
\end{equation}
dove $\lambda_1, \ldots, \lambda_n$ sono gli autovalori di $A$.
\end{definition}

%------------------------------------------------------------
\section{Norme vettoriali}
\label{sec:norme_vettoriali}
%------------------------------------------------------------

\begin{definition}[Norma vettoriale]
Una funzione $\|\cdot\| : \R^n \to \R$ è una \emph{norma} se:
\begin{enumerate}
  \item $\|x\| \geq 0$ e $\|x\| = 0 \iff x = 0$ \quad (positività);
  \item $\|\alpha x\| = |\alpha|\,\|x\|$ per ogni $\alpha \in \R$ \quad (omogeneità);
  \item $\|x + y\| \leq \|x\| + \|y\|$ \quad (disuguaglianza triangolare).
\end{enumerate}
\end{definition}

Le norme $p$ più comuni ($p \geq 1$) sono:
\begin{align}
  \|x\|_1   &= \sum_{i=1}^n |x_i|, \label{eq:norma1}\\
  \|x\|_2   &= \sqrt{\sum_{i=1}^n x_i^2}, \label{eq:norma2}\\
  \|x\|_\infty &= \max_{1 \leq i \leq n} |x_i|. \label{eq:normaInf}
\end{align}

\begin{theorem}[Equivalenza delle norme in $\R^n$]
Tutte le norme su $\R^n$ sono equivalenti: per ogni coppia di norme
$\|\cdot\|_\alpha$, $\|\cdot\|_\beta$ esistono costanti $c_1, c_2 > 0$ tali che
\begin{equation}
  c_1\,\|x\|_\alpha \leq \|x\|_\beta \leq c_2\,\|x\|_\alpha \quad \forall x \in \R^n.
\end{equation}
\end{theorem}

\subsection{Disuguaglianza di Cauchy-Schwarz}

\begin{theorem}[Cauchy-Schwarz]
Per ogni $x, y \in \R^n$:
\begin{equation}
  |x^\top y| \leq \|x\|_2\,\|y\|_2.
  \label{eq:cauchy_schwarz}
\end{equation}
\end{theorem}

%------------------------------------------------------------
\section{Norme matriciali}
\label{sec:norme_matriciali}
%------------------------------------------------------------

\begin{definition}[Norma matriciale indotta]
La norma matriciale \emph{indotta} (o subordinata) dalla norma vettoriale $\|\cdot\|$
è:
\begin{equation}
  \|A\| = \sup_{x \neq 0} \frac{\|Ax\|}{\|x\|} = \max_{\|x\|=1} \|Ax\|.
  \label{eq:norma_mat}
\end{equation}
\end{definition}

Proprietà fondamentale: $\|Ax\| \leq \|A\|\,\|x\|$ (compatibilità).

Le norme indotte più usate per $A \in \R^{m \times n}$:
\begin{align}
  \|A\|_1      &= \max_{1 \leq j \leq n} \sum_{i=1}^m |a_{ij}|
                  \quad \text{(massima somma di colonna)}, \\
  \|A\|_\infty &= \max_{1 \leq i \leq m} \sum_{j=1}^n |a_{ij}|
                  \quad \text{(massima somma di riga)}, \\
  \|A\|_2      &= \sigma_{\max}(A)
                  \quad \text{(massimo valore singolare)}.
\end{align}

\begin{proposition}
Per ogni norma matriciale indotta vale:
\begin{equation}
  \rho(A) \leq \|A\|.
\end{equation}
\end{proposition}

%------------------------------------------------------------
\section{Numero di condizionamento}
\label{sec:condizionamento_mat}
%------------------------------------------------------------

\begin{definition}[Numero di condizionamento]
Per una matrice non singolare $A$, il \emph{numero di condizionamento} rispetto alla
norma $\|\cdot\|$ è:
\begin{equation}
  \kappa(A) = \|A\|\,\|A^{-1}\|.
  \label{eq:cond}
\end{equation}
\end{definition}

\begin{remark}
Poiché $1 = \|I\| = \|AA^{-1}\| \leq \|A\|\,\|A^{-1}\|$, si ha sempre
$\kappa(A) \geq 1$.
\begin{itemize}
  \item $\kappa(A) \approx 1$: matrice \emph{ben condizionata}.
  \item $\kappa(A) \gg 1$: matrice \emph{mal condizionata}.
\end{itemize}
\end{remark}

Il significato pratico di $\kappa(A)$ verrà approfondito nel
\cref{ch:condizionamento}: in sintesi, $\kappa(A)$ indica di quante cifre
significative si può perdere nella soluzione di $Ax = b$ a causa degli errori di
arrotondamento.

%=============================================================
% Capitolo 3 — Metodi Diretti per Sistemi Lineari
%=============================================================
\chapter{Metodi Diretti per Sistemi Lineari}
\label{ch:diretti}

\section{Panoramica}

I \textbf{metodi diretti} risolvono il sistema $Ax = b$ in un numero finito di
operazioni aritmetiche, producendo (in aritmetica esatta) la soluzione esatta.
Il costo è tipicamente $\mathcal{O}(n^3)$ operazioni in virgola mobile.

L'idea centrale è \textbf{fattorizzare} la matrice $A$ nel prodotto di matrici più
semplici — tipicamente triangolari o diagonali — per ridursi a sistemi facilmente
risolvibili.

%------------------------------------------------------------
\section{Sistemi triangolari}
\label{sec:triangolari}
%------------------------------------------------------------

\subsection{Sistema triangolare inferiore}

Un sistema $Lx = b$ con $L$ triangolare inferiore ($l_{ij} = 0$ per $j > i$) si
risolve con la \textbf{sostituzione in avanti} (\emph{forward substitution}):
\begin{equation}
  x_i = \frac{1}{l_{ii}}\left(b_i - \sum_{j=1}^{i-1} l_{ij}\, x_j\right),
  \quad i = 1, 2, \ldots, n.
  \label{eq:fwd_sub}
\end{equation}

\subsection{Sistema triangolare superiore}

Un sistema $Ux = b$ con $U$ triangolare superiore ($u_{ij} = 0$ per $j < i$) si
risolve con la \textbf{sostituzione all'indietro} (\emph{back substitution}):
\begin{equation}
  x_i = \frac{1}{u_{ii}}\left(b_i - \sum_{j=i+1}^{n} u_{ij}\, x_j\right),
  \quad i = n, n-1, \ldots, 1.
  \label{eq:bck_sub}
\end{equation}

Entrambe le sostituzioni hanno costo $\mathcal{O}(n^2)$.

%------------------------------------------------------------
\section{Fattorizzazione LU}
\label{sec:lu}
%------------------------------------------------------------

\begin{definition}[Fattorizzazione LU]
Una \emph{fattorizzazione LU} di $A \in \R^{n \times n}$ è la scrittura
\begin{equation}
  A = LU,
  \label{eq:LU}
\end{equation}
dove $L$ è triangolare inferiore con elementi diagonali unitari ($l_{ii} = 1$) e
$U$ è triangolare superiore.
\end{definition}

Con la fattorizzazione LU il sistema $Ax = b$ si risolve in due passi:
\begin{align}
  Ly &= b \quad \text{(forward substitution)}, \label{eq:forw}\\
  Ux &= y \quad \text{(back substitution)}.  \label{eq:back}
\end{align}
Il costo totale è $\mathcal{O}(n^3)$ per la fattorizzazione e $\mathcal{O}(n^2)$ per
ciascuna risoluzione successiva (utile quando si deve risolvere più sistemi con la
stessa matrice $A$).

%------------------------------------------------------------
\section{Eliminazione di Gauss}
\label{sec:gauss}
%------------------------------------------------------------

L'\textbf{eliminazione di Gauss} è l'algoritmo classico per calcolare la
fattorizzazione LU. A ogni passo $k$ si elimina l'incognita $x_k$ dalle righe
$k+1, \ldots, n$ mediante opportune combinazioni lineari.

\subsection{Un passo dell'eliminazione}

Al passo $k$ si ha a disposizione la matrice parzialmente eliminata
$A^{(k)} = (a^{(k)}_{ij})$. Per ogni riga $i > k$ si calcola il
\emph{moltiplicatore}:
\begin{equation}
  m_{ik} = \frac{a^{(k)}_{ik}}{a^{(k)}_{kk}},
  \label{eq:moltiplicatore}
\end{equation}
e si aggiorna:
\begin{equation}
  a^{(k+1)}_{ij} = a^{(k)}_{ij} - m_{ik}\, a^{(k)}_{kj},
  \quad j = k, k+1, \ldots, n.
  \label{eq:update_gauss}
\end{equation}

L'elemento $a^{(k)}_{kk}$ si chiama \textbf{pivot} al passo $k$.

\begin{algorithm}[htbp]
  \caption{Eliminazione di Gauss (fattorizzazione LU)}
  \label{alg:gauss}
  \begin{algorithmic}[1]
    \Require $A \in \R^{n \times n}$ non singolare
    \Ensure $L, U$ tali che $A = LU$
    \State $U \gets A$; \; $L \gets I_n$
    \For{$k = 1$ \textbf{to} $n-1$}
      \For{$i = k+1$ \textbf{to} $n$}
        \State $m_{ik} \gets u_{ik} / u_{kk}$
        \State $l_{ik} \gets m_{ik}$
        \For{$j = k$ \textbf{to} $n$}
          \State $u_{ij} \gets u_{ij} - m_{ik} \cdot u_{kj}$
        \EndFor
      \EndFor
    \EndFor
    \State \Return $L$, $U$
  \end{algorithmic}
\end{algorithm}

Il costo dell'eliminazione di Gauss è $\frac{2}{3}n^3 + \mathcal{O}(n^2)$
operazioni in virgola mobile.

\begin{theorem}[Esistenza della fattorizzazione LU]
La fattorizzazione LU senza pivoting esiste ed è unica se e solo se tutti i
\emph{minori principali} di $A$ sono non nulli.
\end{theorem}

%------------------------------------------------------------
\section{Pivoting}
\label{sec:pivoting}
%------------------------------------------------------------

Nella pratica il pivot $a^{(k)}_{kk}$ può essere zero (impossibilità) o molto
piccolo (instabilità numerica). Per evitare questi problemi si adotta il
\textbf{pivoting parziale per colonne}: ad ogni passo $k$ si sceglie come pivot
l'elemento di modulo massimo nella colonna $k$ tra le righe $k, k+1, \ldots, n$,
e si scambiano le righe corrispondentemente.

Con il pivoting parziale si ha la fattorizzazione:
\begin{equation}
  PA = LU,
  \label{eq:PLU}
\end{equation}
dove $P$ è una \textbf{matrice di permutazione} che registra gli scambi di riga.

\begin{remark}[Stabilità]
L'eliminazione di Gauss con pivoting parziale è numericamente stabile per quasi
tutte le matrici in pratica, sebbene esistano esempi patologici di instabilità.
Il pivoting totale (su righe e colonne) è più stabile ma meno usato per il
maggiore costo computazionale.
\end{remark}

%------------------------------------------------------------
\section{Matrici sparse e metodi a basso costo}
\label{sec:sparse}
%------------------------------------------------------------

Quando $A$ ha molti elementi nulli (\emph{matrice sparsa}), la fattorizzazione LU
standard non sfrutta la struttura e può produrre un \emph{fill-in} (riempimento
degli zeri) costoso. Tecniche specializzate (reordering, fattorizzazioni incomplete)
permettono di gestire efficientemente sistemi sparsi di grandi dimensioni.

%------------------------------------------------------------
\section{Costo computazionale e confronto}
\label{sec:costo_diretti}
%------------------------------------------------------------

\begin{table}[htbp]
  \centering
  \caption{Costo computazionale dei principali metodi diretti.}
  \label{tab:costo}
  \begin{tabular}{lcc}
    \toprule
    \textbf{Operazione} & \textbf{Flop (ordine)} & \textbf{Note} \\
    \midrule
    Fattorizzazione LU             & $\tfrac{2}{3}n^3$ & una sola volta per $A$ \\
    Forward/back substitution      & $n^2$             & per ogni $b$ \\
    Calcolo esplicito di $A^{-1}$  & $2n^3$            & sconsigliato \\
    \bottomrule
  \end{tabular}
\end{table}

%=============================================================
% Capitolo 4 — Metodi Iterativi per Sistemi Lineari
%=============================================================
\chapter{Metodi Iterativi per Sistemi Lineari}
\label{ch:iterativi}

\section{Motivazione}

I metodi diretti richiedono $\mathcal{O}(n^3)$ operazioni e $\mathcal{O}(n^2)$ di
memoria: per sistemi di grandi dimensioni (e.g.\ $n > 10^6$, tipici in data science e
simulazione) questo è proibitivo. I \textbf{metodi iterativi} costruiscono una
successione $\{x^{(k)}\}$ che converge alla soluzione $x^*$ di $Ax = b$, e spesso
richiedono solo $\mathcal{O}(n)$ operazioni per iterazione se $A$ è sparsa.

%------------------------------------------------------------
\section{Framework generale: splitting della matrice}
\label{sec:splitting}
%------------------------------------------------------------

L'idea comune è \textbf{splittare} la matrice $A$ come:
\begin{equation}
  A = M - N,
  \label{eq:splitting}
\end{equation}
dove $M$ è una matrice scelta in modo che $Mx = c$ sia facile da risolvere.
Sostituendo in $Ax = b$:
\begin{equation}
  Mx = Nx + b \implies x = M^{-1}Nx + M^{-1}b.
  \label{eq:iter_gen}
\end{equation}

Questo definisce la \textbf{iterazione stazionaria}:
\begin{equation}
  x^{(k+1)} = \underbrace{M^{-1}N}_{B}\, x^{(k)} + \underbrace{M^{-1}b}_{c},
  \label{eq:iter_staz}
\end{equation}
dove $B = M^{-1}N = M^{-1}(M - A) = I - M^{-1}A$ è la \textbf{matrice di iterazione}.

Decomponiamo $A$ nei suoi componenti diagonale, triangolare inferiore e superiore:
\begin{equation}
  A = D + L + U,
  \label{eq:DLU}
\end{equation}
dove:
\begin{itemize}
  \item $D = \operatorname{diag}(a_{11}, \ldots, a_{nn})$ è la parte diagonale;
  \item $L$ è la parte strettamente triangolare inferiore ($l_{ij} = a_{ij}$ per $i > j$);
  \item $U$ è la parte strettamente triangolare superiore ($u_{ij} = a_{ij}$ per $i < j$).
\end{itemize}

%------------------------------------------------------------
\section{Metodo di Jacobi}
\label{sec:jacobi}
%------------------------------------------------------------

\begin{definition}[Metodo di Jacobi]
Lo splitting di Jacobi è $M = D$, $N = -(L + U)$. L'iterazione è:
\begin{equation}
  x^{(k+1)} = D^{-1}\bigl(-(L+U)x^{(k)} + b\bigr),
  \label{eq:jacobi_mat}
\end{equation}
ovvero, per ogni componente $i$:
\begin{equation}
  x_i^{(k+1)} = \frac{1}{a_{ii}}\left(b_i - \sum_{\substack{j=1 \\ j \neq i}}^n
  a_{ij}\, x_j^{(k)}\right), \quad i = 1, \ldots, n.
  \label{eq:jacobi_comp}
\end{equation}
\end{definition}

\begin{remark}
Ogni componente $x_i^{(k+1)}$ si calcola usando \emph{solo} le componenti
dell'iterata precedente $x^{(k)}$: il metodo è \textbf{parallelizzabile}.
\end{remark}

La matrice di iterazione di Jacobi è:
\begin{equation}
  B_J = -D^{-1}(L + U) = I - D^{-1}A.
  \label{eq:B_jacobi}
\end{equation}

%------------------------------------------------------------
\section{Metodo di Gauss-Seidel}
\label{sec:gauss_seidel}
%------------------------------------------------------------

\begin{definition}[Metodo di Gauss-Seidel]
Lo splitting di Gauss-Seidel è $M = D + L$, $N = -U$. L'iterazione è:
\begin{equation}
  (D + L)\,x^{(k+1)} = -U\,x^{(k)} + b,
  \label{eq:gs_mat}
\end{equation}
ovvero, per ogni componente $i$:
\begin{equation}
  x_i^{(k+1)} = \frac{1}{a_{ii}}\left(b_i
  - \sum_{j=1}^{i-1} a_{ij}\, x_j^{(k+1)}
  - \sum_{j=i+1}^{n} a_{ij}\, x_j^{(k)}\right), \quad i = 1, \ldots, n.
  \label{eq:gs_comp}
\end{equation}
\end{definition}

\begin{remark}[Spostamenti successivi]
A differenza di Jacobi, Gauss-Seidel usa immediatamente i valori $x_j^{(k+1)}$ già
calcolati (per $j < i$). Questo accelera tipicamente la convergenza, ma rende il
metodo \textbf{non parallelizzabile} nella forma standard.
\end{remark}

La matrice di iterazione di Gauss-Seidel è:
\begin{equation}
  B_{GS} = -(D + L)^{-1}U.
  \label{eq:B_gs}
\end{equation}

%------------------------------------------------------------
\section{Criteri di convergenza}
\label{sec:convergenza}
%------------------------------------------------------------

\subsection{Condizione necessaria e sufficiente}

\begin{theorem}[Convergenza dei metodi iterativi stazionari]
\label{thm:conv_iter}
La successione $x^{(k+1)} = Bx^{(k)} + c$ converge alla soluzione di $Ax = b$
per ogni $x^{(0)}$ se e solo se:
\begin{equation}
  \rho(B) < 1,
  \label{eq:conv_cond}
\end{equation}
dove $\rho(B)$ è il raggio spettrale della matrice di iterazione $B$.
\end{theorem}

Più piccolo è $\rho(B)$, più rapida è la convergenza.

\subsection{Condizione sufficiente: norma}

\begin{corollary}
Se esiste una norma matriciale $\|\cdot\|$ tale che $\|B\| < 1$, allora il metodo
converge.
\end{corollary}

\subsection{Teorema per matrici a dominanza diagonale}

\begin{definition}[Dominanza diagonale stretta per righe]
$A \in \R^{n \times n}$ è \emph{strettamente a dominanza diagonale per righe} se:
\begin{equation}
  |a_{ii}| > \sum_{\substack{j=1 \\ j \neq i}}^n |a_{ij}|, \quad i = 1, \ldots, n.
  \label{eq:dom_diag}
\end{equation}
\end{definition}

\begin{theorem}[Convergenza con dominanza diagonale]
\label{thm:dom_diag}
Se $A$ è strettamente a dominanza diagonale per righe, allora i metodi di Jacobi e
di Gauss-Seidel convergono entrambi.
\end{theorem}

\begin{proof}[Dimostrazione (Jacobi)]
Per la matrice di iterazione $B_J = -D^{-1}(L+U)$:
\begin{equation}
  \|B_J\|_\infty = \max_i \sum_{j \neq i} \frac{|a_{ij}|}{|a_{ii}|} < 1,
\end{equation}
grazie alla dominanza diagonale stretta. Per il corollario precedente, Jacobi converge.
\end{proof}

\begin{example}[Sistema a dominanza diagonale]
Consideriamo il sistema
\[
  \begin{pmatrix} 10 & -1 & 2 \\ -1 & 11 & -1 \\ 2 & -1 & 10 \end{pmatrix}
  \begin{pmatrix} x_1 \\ x_2 \\ x_3 \end{pmatrix}
  =
  \begin{pmatrix} 6 \\ 25 \\ -11 \end{pmatrix}.
\]
La matrice è strettamente a dominanza diagonale ($10 > 1+2$, $11 > 1+1$, $10 > 2+1$),
quindi Jacobi e Gauss-Seidel convergono entrambi.
\end{example}

%------------------------------------------------------------
\section{Velocità di convergenza e criterio di arresto}
\label{sec:stop}
%------------------------------------------------------------

L'errore al passo $k$ si propaga come:
\begin{equation}
  e^{(k)} = x^{(k)} - x^* = B^k e^{(0)},
\end{equation}
quindi $\|e^{(k)}\| \leq \|B\|^k \|e^{(0)}\|$.

Il criterio di arresto più comune è sul \textbf{residuo relativo}:
\begin{equation}
  \frac{\|b - Ax^{(k)}\|}{\|b\|} < \text{tol},
  \label{eq:stop}
\end{equation}
dove \texttt{tol} è la tolleranza desiderata (e.g.\ $10^{-8}$).

\begin{remark}[Gauss-Seidel vs.\ Jacobi]
In generale $\rho(B_{GS}) \leq \rho(B_J)^2$ (per matrici definite positive
simmetriche): Gauss-Seidel converge circa il doppio più velocemente di Jacobi in
termini di iterazioni. Tuttavia Jacobi è più adatto alla parallelizzazione.
\end{remark}

%=============================================================
% Capitolo 5 — Condizionamento dei Sistemi Lineari
%=============================================================
\chapter{Condizionamento dei Sistemi Lineari}
\label{ch:condizionamento}

\section{Il problema della perturbazione}

Nella pratica i dati di un sistema lineare $Ax = b$ non sono mai noti con precisione
infinita: sono affetti da errori di misura, di arrotondamento, o di approssimazione.
È fondamentale capire \emph{quanto} la soluzione $x$ cambia al variare (anche
leggermente) dei dati $A$ e $b$.

%------------------------------------------------------------
\section{Perturbazione solo sul termine noto}
\label{sec:pert_b}
%------------------------------------------------------------

Consideriamo il sistema perturbato
\begin{equation}
  A(x + \delta x) = b + \delta b,
  \label{eq:pert_b}
\end{equation}
dove $\delta b$ è una piccola perturbazione del termine noto e $\delta x$ è la
conseguente perturbazione della soluzione. Sottraendo $Ax = b$:
\begin{equation}
  A\,\delta x = \delta b \implies \delta x = A^{-1}\delta b.
\end{equation}

Prendendo le norme:
\begin{equation}
  \|\delta x\| \leq \|A^{-1}\|\,\|\delta b\|.
  \label{eq:bound_deltax}
\end{equation}
Poiché $\|b\| = \|Ax\| \leq \|A\|\,\|x\|$, si ottiene $\frac{1}{\|x\|} \leq
\frac{\|A\|}{\|b\|}$. Combinando:

\begin{equation}
  \frac{\|\delta x\|}{\|x\|} \leq \|A\|\,\|A^{-1}\|\,\frac{\|\delta b\|}{\|b\|}
  = \kappa(A)\,\frac{\|\delta b\|}{\|b\|}.
  \label{eq:cond_b}
\end{equation}

%------------------------------------------------------------
\section{Perturbazione sulla matrice e sul termine noto}
\label{sec:pert_Ab}
%------------------------------------------------------------

Consideriamo ora il caso più generale in cui sia $A$ che $b$ sono perturbati:
\begin{equation}
  (A + \delta A)(x + \delta x) = b + \delta b.
  \label{eq:pert_Ab}
\end{equation}

Espandendo e trascurando i termini del secondo ordine ($\delta A\,\delta x \approx 0$):
\begin{equation}
  A\,\delta x \approx \delta b - \delta A\, x,
\end{equation}
da cui $\delta x = A^{-1}(\delta b - \delta A\, x)$, e prendendo le norme:

\begin{theorem}[Limitazione dell'errore relativo — caso generale]
\label{thm:cond_gen}
Se $\|\delta A\| < \|A^{-1}\|^{-1}$ (perturbazione piccola), allora:
\begin{equation}
  \frac{\|\delta x\|}{\|x + \delta x\|} \leq
  \frac{\kappa(A)}{1 - \kappa(A)\,\frac{\|\delta A\|}{\|A\|}}
  \left(\frac{\|\delta A\|}{\|A\|} + \frac{\|\delta b\|}{\|b\|}\right).
  \label{eq:cond_gen}
\end{equation}
\end{theorem}

\begin{remark}
Il \cref{thm:cond_gen} mostra che l'errore relativo sulla soluzione è amplificato dal
fattore $\kappa(A)$: se $\kappa(A) = 10^k$, si perdono fino a $k$ cifre significative
nella soluzione rispetto alla precisione dei dati.
\end{remark}

%------------------------------------------------------------
\section{Interpretazione pratica del condizionamento}
\label{sec:interpretazione}
%------------------------------------------------------------

In aritmetica in virgola mobile con epsilon di macchina $\eps_M$, le perturbazioni
sui dati sono dell'ordine di $\eps_M$. La \cref{eq:cond_gen} implica che l'errore
relativo sulla soluzione è approssimativamente:
\begin{equation}
  \frac{\|\delta x\|}{\|x\|} \approx \kappa(A)\,\eps_M.
  \label{eq:errore_pratico}
\end{equation}

\begin{example}[Sistema mal condizionato]
La \textbf{matrice di Hilbert} di ordine $n$, con elementi $h_{ij} = 1/(i+j-1)$,
è notoriamente mal condizionata:
\[
  H_5 = \begin{pmatrix}
    1 & \frac{1}{2} & \frac{1}{3} & \frac{1}{4} & \frac{1}{5} \\[4pt]
    \frac{1}{2} & \frac{1}{3} & \frac{1}{4} & \frac{1}{5} & \frac{1}{6} \\[4pt]
    \frac{1}{3} & \frac{1}{4} & \frac{1}{5} & \frac{1}{6} & \frac{1}{7} \\[4pt]
    \frac{1}{4} & \frac{1}{5} & \frac{1}{6} & \frac{1}{7} & \frac{1}{8} \\[4pt]
    \frac{1}{5} & \frac{1}{6} & \frac{1}{7} & \frac{1}{8} & \frac{1}{9}
  \end{pmatrix}, \quad \kappa_2(H_5) \approx 4.77 \times 10^5.
\]
Con doppia precisione ($\eps_M \approx 10^{-16}$) l'errore relativo sulla soluzione
è dell'ordine di $4.77 \times 10^5 \times 10^{-16} \approx 10^{-11}$: si perdono
circa 5 cifre significative.
\end{example}

%------------------------------------------------------------
\section{Stima del numero di condizionamento}
\label{sec:stima_cond}
%------------------------------------------------------------

Calcolare $\kappa(A) = \|A\|\,\|A^{-1}\|$ esattamente richiede di conoscere $A^{-1}$,
il che è costoso ($\mathcal{O}(n^3)$). In pratica:
\begin{itemize}
  \item Per la norma 2: $\kappa_2(A) = \sigma_{\max}(A)/\sigma_{\min}(A)$
        (rapporto tra valore singolare massimo e minimo), calcolabile con la SVD.
  \item Metodi di stima a basso costo (e.g.\ \texttt{rcond} in MATLAB/NumPy)
        forniscono una stima di $1/\kappa(A)$ in $\mathcal{O}(n^2)$ una volta
        calcolata la fattorizzazione LU.
\end{itemize}

%------------------------------------------------------------
\section{Residuo e soluzione calcolata}
\label{sec:residuo}
%------------------------------------------------------------

Sia $\hat{x}$ la soluzione calcolata numericamente di $Ax = b$. Il
\textbf{residuo} è:
\begin{equation}
  r = b - A\hat{x}.
  \label{eq:residuo}
\end{equation}

Un residuo piccolo \textbf{non} implica necessariamente un errore piccolo. Il legame
è:
\begin{equation}
  \frac{\|x - \hat{x}\|}{\|x\|} \leq \kappa(A)\,\frac{\|r\|}{\|b\|}.
  \label{eq:residuo_errore}
\end{equation}

\begin{remark}[Backward e forward error]
\begin{itemize}
  \item Il \emph{backward error} (errore all'indietro) è $\|r\|/\|b\|$: misura di
        quanto $\hat{x}$ risolve \emph{esattamente} un sistema perturbato.
  \item Il \emph{forward error} (errore in avanti) è $\|x - \hat{x}\|/\|x\|$:
        misura la distanza dalla soluzione vera.
\end{itemize}
Un metodo si dice \textbf{backward-stabile} se il backward error è dell'ordine di
$\eps_M$.
\end{remark}

%------------------------------------------------------------
\section{Riepilogo}
\label{sec:riepilogo_cond}
%------------------------------------------------------------

\begin{table}[htbp]
  \centering
  \caption{Effetto del numero di condizionamento sulla soluzione (doppia precisione,
           $\eps_M \approx 10^{-16}$).}
  \label{tab:condizionamento}
  \begin{tabular}{ccc}
    \toprule
    $\kappa(A)$ & Cifre perse & Qualità \\
    \midrule
    $\sim 1$       & $\sim 0$  & Eccellente \\
    $\sim 10^3$    & $\sim 3$  & Buona \\
    $\sim 10^8$    & $\sim 8$  & Accettabile \\
    $\sim 10^{12}$ & $\sim 12$ & Scarsa \\
    $\sim 10^{16}$ & $\sim 16$ & Soluzione inattendibile \\
    \bottomrule
  \end{tabular}
\end{table}

In sintesi: prima di risolvere un sistema lineare è \emph{sempre} buona pratica
stimare $\kappa(A)$ per capire quante cifre significative ci si può aspettare nella
soluzione.

%=============================================================
% Capitolo 6 — Interpolazione Polinomiale
%=============================================================
\chapter{Interpolazione e Approssimazione di Funzioni}
\label{ch:interpolazione}

\section{Motivazione}

In molte applicazioni si dispone di un insieme di \emph{dati discreti}
$\{(x_i, f_i)\}_{i=0}^n$ (misure, simulazioni, valori di una funzione ignota) e si
vuole:
\begin{itemize}
  \item \textbf{Interpolare}: trovare una funzione $p$ che passi esattamente per tutti
        i punti, cioè $p(x_i) = f_i$ per $i = 0,\ldots,n$;
  \item \textbf{Approssimare}: trovare una funzione $p$ che si avvicini ai dati nel
        senso di qualche norma, senza necessariamente passare per tutti i punti.
\end{itemize}

%------------------------------------------------------------
\section{Interpolazione polinomiale}
\label{sec:interp_polin}
%------------------------------------------------------------

La scelta più naturale per $p$ è un \textbf{polinomio}: lo spazio
$\mathcal{P}_n$ dei polinomi di grado $\leq n$ ha dimensione $n+1$, quindi con
$n+1$ punti si ha (genericamente) un unico polinomio interpolante.

\begin{theorem}[Esistenza e unicità del polinomio interpolante]
Dati $n+1$ punti distinti $x_0 < x_1 < \cdots < x_n$ e valori $f_0,\ldots,f_n$,
esiste un unico polinomio $p_n \in \mathcal{P}_n$ tale che
\begin{equation}
  p_n(x_i) = f_i, \quad i = 0, 1, \ldots, n.
  \label{eq:interp_cond}
\end{equation}
\end{theorem}

\subsection{Matrice di Vandermonde}

Nella \emph{base delle potenze} $\{1, x, x^2, \ldots, x^n\}$, scrivendo
$p_n(x) = \sum_{j=0}^n a_j x^j$, le condizioni di interpolazione diventano il
sistema lineare:
\begin{equation}
  \underbrace{
  \begin{pmatrix}
    1      & x_0    & x_0^2  & \cdots & x_0^n \\
    1      & x_1    & x_1^2  & \cdots & x_1^n \\
    \vdots & \vdots & \vdots & \ddots & \vdots \\
    1      & x_n    & x_n^2  & \cdots & x_n^n
  \end{pmatrix}
  }_{V \text{ (matrice di Vandermonde)}}
  \begin{pmatrix} a_0 \\ a_1 \\ \vdots \\ a_n \end{pmatrix}
  =
  \begin{pmatrix} f_0 \\ f_1 \\ \vdots \\ f_n \end{pmatrix}.
  \label{eq:vandermonde}
\end{equation}

\begin{remark}
La matrice di Vandermonde è invertibile se e solo se le ascisse $x_i$ sono tutte
distinte. Tuttavia, per $n$ grande $V$ tende ad essere \emph{mal condizionata}: è
preferibile usare la base di Lagrange.
\end{remark}

%------------------------------------------------------------
\section{Polinomio di Lagrange}
\label{sec:lagrange}
%------------------------------------------------------------

\begin{definition}[Polinomi di base di Lagrange]
I \emph{polinomi di base di Lagrange} associati alle ascisse $x_0,\ldots,x_n$ sono:
\begin{equation}
  L_k^{(n)}(x) = \prod_{\substack{j=0 \\ j \neq k}}^n \frac{x - x_j}{x_k - x_j},
  \quad k = 0, 1, \ldots, n.
  \label{eq:lagrange_base}
\end{equation}
\end{definition}

Notare che $L_k^{(n)}(x_j) = \delta_{kj}$ (delta di Kronecker). Il polinomio
interpolante si scrive allora in \emph{forma di Lagrange}:
\begin{equation}
  p_n(x) = \sum_{k=0}^n f_k\, L_k^{(n)}(x).
  \label{eq:lagrange_polin}
\end{equation}

\begin{remark}
La forma di Lagrange mostra che \emph{stesso spazio, stessa soluzione}: qualunque
base si usi per $\mathcal{P}_n$, il polinomio interpolante è unico.
\end{remark}

%------------------------------------------------------------
\section{Errore di interpolazione}
\label{sec:errore_interp}
%------------------------------------------------------------

\begin{theorem}[Errore del polinomio interpolante]
\label{thm:errore_interp}
Sia $f \in C^{n+1}([a,b])$ e $p_n$ il polinomio interpolante nei punti
$x_0,\ldots,x_n \in [a,b]$. Allora, per ogni $\bar{x} \in [a,b]$, esiste
$\xi = \xi(\bar{x}) \in (a,b)$ tale che:
\begin{equation}
  e_n(\bar{x}) = f(\bar{x}) - p_n(\bar{x})
  = \frac{f^{(n+1)}(\xi)}{(n+1)!}\,\omega_{n+1}(\bar{x}),
  \label{eq:errore_interp}
\end{equation}
dove
\begin{equation}
  \omega_{n+1}(x) = \prod_{i=0}^n (x - x_i).
  \label{eq:omega}
\end{equation}
\end{theorem}

\begin{proof}[Schema della dimostrazione]
Si definisce $g(x) = f(x) - p_n(x) - R_n(\bar{x})\,\omega_{n+1}(x)$, dove
$R_n(\bar{x})$ è scelto affinché $g(\bar{x}) = 0$. Allora $g$ ha (almeno) $n+2$
zeri in $[a,b]$: i punti $x_0,\ldots,x_n$ e $\bar{x}$. Per il teorema di Rolle,
$g^{(n+1)}$ ha almeno uno zero $\xi$ in $(a,b)$, da cui si ricava la formula.
\end{proof}

\subsection{Stima dell'errore}

Maggiorando \cref{eq:errore_interp}:
\begin{equation}
  \|e_n\|_\infty \leq \frac{M_{n+1}}{(n+1)!}\,\|\omega_{n+1}\|_\infty,
  \label{eq:bound_errore_interp}
\end{equation}
dove $M_{n+1} = \max_{x \in [a,b]} |f^{(n+1)}(x)|$ e
$\|\omega_{n+1}\|_\infty = \max_{x \in [a,b]} |\omega_{n+1}(x)|$.

\begin{example}
Per nodi equispaziati su $[a,b]$ con passo $h = (b-a)/n$:
\[
  \|\omega_{n+1}\|_\infty \leq \frac{h^{n+1}\, n!}{4},
\]
quindi $\|e_n\|_\infty \leq \frac{M_{n+1}\,h^{n+1}}{4(n+1)}$.
\end{example}

%------------------------------------------------------------
\section{Costante di Lebesgue}
\label{sec:lebesgue}
%------------------------------------------------------------

\begin{definition}[Costante di Lebesgue]
La \emph{costante di Lebesgue} associata alle ascisse $x_0,\ldots,x_n$ è:
\begin{equation}
  \Lambda_n = \max_{x \in [a,b]} \sum_{k=0}^n |L_k^{(n)}(x)|.
  \label{eq:lebesgue}
\end{equation}
\end{definition}

La costante di Lebesgue misura il \emph{condizionamento} del problema di
interpolazione:

\begin{theorem}[Legame tra errore di interpolazione e costante di Lebesgue]
Sia $p^*_n$ la migliore approssimazione polinomiale di $f$ in $\mathcal{P}_n$
(nel senso della norma $\|\cdot\|_\infty$). Allora:
\begin{equation}
  \|f - p_n\|_\infty \leq (1 + \Lambda_n)\,\|f - p^*_n\|_\infty.
  \label{eq:lebesgue_bound}
\end{equation}
\end{theorem}

\begin{remark}[Crescita di $\Lambda_n$]
Per nodi equispaziati, $\Lambda_n$ cresce esponenzialmente con $n$
(\emph{fenomeno di Runge}): l'interpolazione polinomiale di alto grado con nodi
equispaziati può divergere. Nodi di Chebyshev minimizzano $\|\omega_{n+1}\|_\infty$
e producono $\Lambda_n = \mathcal{O}(\log n)$.
\end{remark}

\begin{remark}[Spline]
Per evitare il fenomeno di Runge, in pratica si preferisce usare le
\textbf{funzioni spline}: polinomi di basso grado (tipicamente cubiche) raccordati
con continuità sui sotto-intervalli. L'errore converge a zero al crescere di $n$
senza fenomeni di oscillazione.
\end{remark}

%=============================================================
% Capitolo 7 — Minimi Quadrati e Fattorizzazione QR
%=============================================================
\chapter{Minimi Quadrati e Fattorizzazione QR}
\label{ch:minimi_quadrati}

\section{Il problema dei minimi quadrati}
\label{sec:mq_problema}

Quando il numero di dati $m$ supera il numero di parametri $n$ (sistema
\emph{sovradeterminato}), in generale non esiste $x \in \R^n$ tale che $Ax = b$
con $A \in \R^{m \times n}$, $m > n$. Il problema dei \textbf{minimi quadrati}
chiede di minimizzare il \emph{residuo}:
\begin{equation}
  \min_{x \in \R^n} \|Ax - b\|_2^2 = \min_{x \in \R^n} \sum_{i=1}^m \bigl(
  (Ax)_i - b_i \bigr)^2.
  \label{eq:mq}
\end{equation}

\subsection{Applicazione: fitting di dati}

Dati $m$ punti $\{(x_i, f_i)\}_{i=1}^m$, si cerca una funzione
$y(x) = \sum_{j=1}^n \alpha_j \varphi_j(x)$ (con $\varphi_j$ funzioni base
fissate) che minimizzi:
\begin{equation}
  S = \sum_{i=1}^m \bigl(f_i - y(x_i)\bigr)^2.
  \label{eq:mq_fit}
\end{equation}
Questo è esattamente \cref{eq:mq} con $a_{ij} = \varphi_j(x_i)$.

%------------------------------------------------------------
\section{Equazioni normali}
\label{sec:equazioni_normali}
%------------------------------------------------------------

La soluzione del problema \cref{eq:mq} soddisfa le \textbf{equazioni normali}:
\begin{equation}
  A^\top A\, x = A^\top b.
  \label{eq:normali}
\end{equation}

Se $A$ ha rango pieno ($\operatorname{rank}(A) = n$), la matrice $A^\top A$ è
definita positiva e la soluzione è unica:
\begin{equation}
  x = (A^\top A)^{-1} A^\top b.
  \label{eq:sol_mq}
\end{equation}

\begin{remark}[Mal condizionamento]
La matrice $A^\top A$ ha numero di condizionamento $\kappa_2(A^\top A) =
\kappa_2(A)^2$: risolvere le equazioni normali raddoppia il numero di cifre
perse. È preferibile usare la fattorizzazione QR.
\end{remark}

%------------------------------------------------------------
\section{Fattorizzazione QR}
\label{sec:qr}
%------------------------------------------------------------

\begin{definition}[Fattorizzazione QR]
Per $A \in \R^{m \times n}$ con $m \geq n$ e $\operatorname{rank}(A) = n$, la
\emph{fattorizzazione QR} è:
\begin{equation}
  A = QR = Q \begin{pmatrix} \hat{R} \\ 0 \end{pmatrix},
  \label{eq:QR}
\end{equation}
dove $Q \in \R^{m \times m}$ è ortogonale ($Q^\top Q = I$) e
$\hat{R} \in \R^{n \times n}$ è triangolare superiore con elementi diagonali
positivi (se si richiede unicità).
\end{definition}

\subsection{Soluzione del problema dei minimi quadrati tramite QR}

Poiché la moltiplicazione per una matrice ortogonale non cambia la norma
euclidea:
\begin{align}
  \|Ax - b\|_2^2 &= \|QRx - b\|_2^2 = \|Q(Rx - Q^\top b)\|_2^2
                  = \|Rx - Q^\top b\|_2^2 \nonumber \\
  &= \left\|\begin{pmatrix}\hat{R}x - g_1 \\ -g_2\end{pmatrix}\right\|_2^2
   = \|\hat{R}x - g_1\|_2^2 + \|g_2\|_2^2,
  \label{eq:mq_qr}
\end{align}
dove $Q^\top b = \begin{pmatrix}g_1 \\ g_2\end{pmatrix}$ con
$g_1 \in \R^n$, $g_2 \in \R^{m-n}$.

Il minimo si raggiunge quando $\hat{R}x = g_1$, sistema triangolare
superiore facilmente risolubile. Il residuo minimo vale $\|g_2\|_2$.

%------------------------------------------------------------
\section{Matrici di Householder}
\label{sec:householder}
%------------------------------------------------------------

Le \textbf{riflessioni di Householder} permettono di calcolare la fattorizzazione
QR in modo numericamente stabile.

\begin{definition}[Matrice di Householder]
Data una direzione $v \in \R^m$ con $\|v\|_2 = 1$, la \emph{matrice di
Householder} è:
\begin{equation}
  H = I - 2vv^\top.
  \label{eq:householder}
\end{equation}
\end{definition}

$H$ è \emph{simmetrica} e \emph{ortogonale}: $H^\top = H$ e $H^2 = I$, quindi
$H^{-1} = H^\top = H$.

\subsection{Costruzione della fattorizzazione QR}

Dato un vettore $z \in \R^m$, si sceglie $v$ tale che $Hz = \pm\|z\|_2\,e_1$.
Applicando successivamente $n$ riflessioni:
\begin{equation}
  H_n \cdots H_2 H_1\, A = \begin{pmatrix}\hat{R} \\ 0\end{pmatrix},
  \label{eq:QR_householder}
\end{equation}
da cui $Q^\top = H_n \cdots H_1$ e $Q = H_1 \cdots H_n$ (poiché $H_k^{-1} = H_k$).

\begin{algorithm}[htbp]
  \caption{Fattorizzazione QR tramite riflessioni di Householder}
  \label{alg:qr_householder}
  \begin{algorithmic}[1]
    \Require $A \in \R^{m \times n}$ con $m \geq n$, $\operatorname{rank}(A)=n$
    \Ensure $Q \in \R^{m \times m}$ ortogonale, $\hat{R} \in \R^{n \times n}$
            triangolare superiore, $A = Q\begin{pmatrix}\hat{R}\\0\end{pmatrix}$
    \State $R \gets A$; \; $Q \gets I_m$
    \For{$k = 1$ \textbf{to} $n$}
      \State $z \gets R(k:m, k)$ \Comment{colonna $k$ dal pivot in giù}
      \State $v \gets z + \operatorname{sign}(z_1)\|z\|_2\,e_1$; \;
             $v \gets v / \|v\|_2$
      \State $H_k \gets I_{m-k+1} - 2vv^\top$
      \State $R(k:m, k:n) \gets H_k\, R(k:m, k:n)$
      \State $Q(1:m, k:m) \gets Q(1:m, k:m)\, H_k^\top$
    \EndFor
    \State \Return $Q$, $\hat{R} \gets R(1:n, 1:n)$
  \end{algorithmic}
\end{algorithm}

\subsection{Confronto: QR vs.\ equazioni normali}

\begin{table}[htbp]
  \centering
  \caption{Confronto tra i metodi per i minimi quadrati.}
  \label{tab:mq_confronto}
  \begin{tabular}{lcc}
    \toprule
    \textbf{Metodo} & \textbf{Costo} & \textbf{Stabilità} \\
    \midrule
    Equazioni normali ($A^\top A$) & $mn^2 + \frac{n^3}{3}$ & Scarsa \\
    Fattorizzazione QR (Householder) & $2mn^2 - \frac{2n^3}{3}$ & Buona \\
    SVD                              & $\mathcal{O}(mn^2)$ & Ottima (usata quando $A$ è quasi singolare) \\
    \bottomrule
  \end{tabular}
\end{table}

%=============================================================
% Capitolo 8 — Autovalori e Decomposizione SVD
%=============================================================
\chapter{Autovalori e Decomposizione ai Valori Singolari}
\label{ch:autovalori}

\section{Autovalori e autovettori}
\label{sec:autovalori}

\begin{definition}
Sia $A \in \mathbb{C}^{n \times n}$. Un numero $\lambda \in \mathbb{C}$ è un
\emph{autovalore} di $A$ e $v \in \mathbb{C}^n \setminus \{0\}$ il corrispondente
\emph{autovettore} se
\begin{equation}
  Av = \lambda v.
  \label{eq:autovettore}
\end{equation}
\end{definition}

\subsection{Equazione caratteristica}

$\lambda$ è autovalore di $A$ se e solo se $\det(A - \lambda I) = 0$.
Il polinomio $p(\lambda) = \det(A - \lambda I)$ di grado $n$ è detto
\textbf{polinomio caratteristico}.

\begin{remark}
In generale, il calcolo degli autovalori tramite il polinomio caratteristico è
\emph{numericamente instabile} per $n \geq 5$ (teorema di Abel-Ruffini). In pratica
si utilizzano metodi iterativi (ad esempio l'algoritmo QR).
\end{remark}

\subsection{Proprietà fondamentali}

Per $A \in \mathbb{C}^{n \times n}$ con autovalori $\lambda_1, \ldots, \lambda_n$:
\begin{align}
  \det(A)  &= \prod_{i=1}^n \lambda_i, \label{eq:det_autoval}\\
  \operatorname{tr}(A) &= \sum_{i=1}^n \lambda_i. \label{eq:traccia}
\end{align}

\subsection{Molteplicità algebrica e geometrica}

\begin{definition}
\begin{itemize}
  \item La \emph{molteplicità algebrica} $\mu_a(\lambda)$ è la molteplicità di
        $\lambda$ come radice del polinomio caratteristico.
  \item La \emph{molteplicità geometrica} $\mu_g(\lambda) = \dim\ker(A-\lambda I)$
        è il numero di autovettori linearmente indipendenti.
\end{itemize}
Si ha sempre $\mu_g(\lambda) \leq \mu_a(\lambda)$.
\end{definition}

%------------------------------------------------------------
\section{Raggio spettrale e cerchi di Gershgorin}
\label{sec:gershgorin}
%------------------------------------------------------------

\begin{theorem}[Cerchi di Gershgorin]
\label{thm:gershgorin}
Ogni autovalore di $A$ appartiene all'unione dei \emph{cerchi di Gershgorin}:
\begin{equation}
  \bigcup_{i=1}^n \mathcal{D}_i, \quad \text{con } \mathcal{D}_i =
  \left\{ z \in \mathbb{C} : |z - a_{ii}| \leq r_i \right\},\quad
  r_i = \sum_{j \neq i} |a_{ij}|.
  \label{eq:gershgorin}
\end{equation}
\end{theorem}

I cerchi di Gershgorin forniscono una \emph{stima rapida e grossolana} della
localizzazione degli autovalori, utile come sanity check prima di calcoli più
accurati.

%------------------------------------------------------------
\section{Diagonalizzazione}
\label{sec:diagonalizzazione}
%------------------------------------------------------------

\begin{definition}[Matrice diagonalizzabile]
$A \in \mathbb{C}^{n \times n}$ è \emph{diagonalizzabile} se esiste una matrice
invertibile $V$ tale che
\begin{equation}
  A = V \Lambda V^{-1}, \quad \Lambda = \operatorname{diag}(\lambda_1, \ldots, \lambda_n).
  \label{eq:diag}
\end{equation}
Le colonne di $V$ sono gli autovettori di $A$.
\end{definition}

\begin{theorem}[Matrici hermitiane]
Se $A = A^*$ (hermitiana; reale: $A = A^\top$), allora:
\begin{enumerate}
  \item Tutti gli autovalori sono \emph{reali};
  \item Gli autovettori relativi ad autovalori distinti sono \emph{ortogonali};
  \item $A$ è \emph{unitariamente diagonalizzabile}: $A = Q\Lambda Q^*$ con $Q$
        unitaria.
\end{enumerate}
\end{theorem}

%------------------------------------------------------------
\section{Decomposizione ai valori singolari (SVD)}
\label{sec:svd}
%------------------------------------------------------------

La SVD estende la diagonalizzazione a matrici \emph{rettangolari}.

\begin{theorem}[Decomposizione ai valori singolari]
\label{thm:svd}
Per ogni $A \in \R^{m \times n}$ esiste la fattorizzazione
\begin{equation}
  A = U \Sigma V^\top,
  \label{eq:svd}
\end{equation}
dove:
\begin{itemize}
  \item $U \in \R^{m \times m}$ è ortogonale (le colonne $u_i$ sono i
        \emph{vettori singolari sinistri});
  \item $V \in \R^{n \times n}$ è ortogonale (le colonne $v_i$ sono i
        \emph{vettori singolari destri});
  \item $\Sigma = \operatorname{diag}(\sigma_1, \ldots, \sigma_p) \in \R^{m \times n}$
        con $p = \min(m,n)$ e $\sigma_1 \geq \sigma_2 \geq \cdots \geq \sigma_p \geq 0$
        (\emph{valori singolari}).
\end{itemize}
\end{theorem}

\subsection{Relazione tra SVD e autovalori}

I valori singolari di $A$ sono le radici quadrate degli autovalori non negativi
di $A^\top A$ (e di $AA^\top$):
\begin{equation}
  \sigma_i = \sqrt{\lambda_i(A^\top A)} = \sqrt{\lambda_i(AA^\top)}.
  \label{eq:svd_autoval}
\end{equation}
Se $A$ è simmetrica definita positiva, $\sigma_i = \lambda_i$.

\subsection{SVD troncata e approssimazione di rango basso}

\begin{theorem}[Teorema di Eckart-Young]
La migliore approssimazione di rango $k$ di $A$ nella norma spettrale e nella
norma di Frobenius è:
\begin{equation}
  A_k = \sum_{i=1}^k \sigma_i\, u_i v_i^\top,
  \label{eq:svd_troncata}
\end{equation}
con errore $\|A - A_k\|_2 = \sigma_{k+1}$.
\end{theorem}

\subsection{Applicazioni della SVD}

\begin{itemize}
  \item \textbf{Norma spettrale}: $\|A\|_2 = \sigma_1$.
  \item \textbf{Numero di condizionamento}: $\kappa_2(A) = \sigma_1 / \sigma_n$ (per
        $A$ quadrata non singolare).
  \item \textbf{Rango numerico}: $\operatorname{rank}(A) = $ numero di valori singolari
        significativamente non nulli.
  \item \textbf{Pseudoinversa di Moore-Penrose}: $A^+ = V\Sigma^+ U^\top$ dove
        $\Sigma^+$ inverte i valori singolari non nulli.
  \item \textbf{Compressione dati}: la SVD troncata è usata in PCA, compressione di
        immagini, sistemi di raccomandazione.
\end{itemize}

%------------------------------------------------------------
\section{Algoritmo QR per il calcolo degli autovalori}
\label{sec:algo_qr}
%------------------------------------------------------------

L'\textbf{algoritmo QR} è il metodo standard per calcolare tutti gli autovalori
di una matrice:

\begin{algorithm}[htbp]
  \caption{Algoritmo QR (forma base)}
  \label{alg:qr_autoval}
  \begin{algorithmic}[1]
    \Require $A_0 = A \in \R^{n \times n}$, numero iterazioni $N$
    \Ensure Approssimazione degli autovalori di $A$
    \For{$k = 0, 1, 2, \ldots, N$}
      \State Calcola la fattorizzazione $A_k = Q_k R_k$
      \State $A_{k+1} \gets R_k Q_k$
    \EndFor
    \State \Return gli elementi diagonali di $A_N$
  \end{algorithmic}
\end{algorithm}

La sequenza $A_k$ converge (sotto opportune ipotesi) a una matrice triangolare
superiore (o quasi-triangolare reale) con gli autovalori sulla diagonale.

In pratica, prima si riduce $A$ a forma di Hessenberg superiore
(costo $\mathcal{O}(n^3)$), poi si applica l'algoritmo QR alla forma ridotta
(ogni iterazione costa $\mathcal{O}(n^2)$).


%------------------------------------------------------------
% Bibliografia
%------------------------------------------------------------
\bibliographystyle{plain}
\begin{thebibliography}{9}

\bibitem{quarteroni}
A.\ Quarteroni, R.\ Sacco, F.\ Saleri,
\textit{Matematica Numerica},
4a ed., Springer, 2014.

\bibitem{nocedal}
J.\ Nocedal, S.\ J.\ Wright,
\textit{Numerical Optimization},
2nd ed., Springer, 2006.

\bibitem{trefethen}
L.\ N.\ Trefethen, D.\ Bau III,
\textit{Numerical Linear Algebra},
SIAM, 1997.

\bibitem{ieee754}
\textit{IEEE Standard for Floating-Point Arithmetic},
IEEE Std 754-2019, IEEE, 2019.

\end{thebibliography}

\end{document}
